\chapter{Technické ukázky}
\label{chap:examples}

Tato kapitola předvádí základní prvky sazby. Při psaní ji nahraďte vlastním
textem. Kniha může být citována samostatně \cite{latex-companion} a později
znovu \cite{latex-companion}. Webová dokumentace je uvedena v~\cite{latex-docs}.
Více zdrojů lze odkázat jedním příkazem \cite{knuth1984,grobid2009}.

\section{Odstavec a seznam}

Odstavce oddělujte prázdným řádkem. Krátký výčet lze zapsat například takto:
\begin{itemize}
  \item první položka,
  \item druhá položka,
  \item poslední položka.
\end{itemize}

\section{Rovnice a odkaz}

Rovnice~\eqref{eq:example} má vlastní štítek, na který lze odkazovat z~textu.
\begin{equation}
  a^2 + b^2 = c^2
  \label{eq:example}
\end{equation}

\section{Tabulka}

Tabulka~\ref{tab:example} ukazuje popisek a odkaz.
\begin{table}[ht]
  \centering
  \caption{Ukázková struktura tabulky}
  \label{tab:example}
  \begin{tabular}{ll}
    \toprule
    Sloupec A & Sloupec B \\
    \midrule
    položka 1 & hodnota 1 \\
    položka 2 & hodnota 2 \\
    \bottomrule
  \end{tabular}
\end{table}

\section{Obrázek}

Obrázek~\ref{fig:example} používá pouze rámeček, aby šablona nepotřebovala
další ukázkový binární soubor.
\begin{figure}[ht]
  \centering
  \fbox{\rule{0pt}{35mm}\rule{0.65\textwidth}{0pt}}
  \caption{Místo pro ukázkový obrázek}
  \label{fig:example}
\end{figure}

% Vlastní soubor vložte do figures/ a rámeček výše nahraďte například:
% \includegraphics[width=0.8\textwidth]{nazev-obrazku.pdf}
